\clearpage
\onecolumn
\appendix
\section{Further proofs and technical details}
We provide additional proofs, detailed versions of the test constructions, and the experimental protocol. The definitions and observation assumptions are those stated in the main text unless explicitly strengthened below.

\section{Localization and Statistical Separation}
\label{arch:sec:localization}

\begin{proposition}[Localization of new violations]
\label{arch:prop:localization}
If $P_0$ is fair on $G$ and $\Impact$ is distributionally sound, every fairness
violation newly introduced by $P_1$ lies in $\ImpEdges$.
\end{proposition}

\begin{proof}
For $e=(u,v)\notin\ImpEdges$, both endpoints lie outside $\Impact$.  Soundness
implies that both endpoint output distributions, and therefore both expectations,
are unchanged.  Hence $g_{1,e}=g_{0,e}$, and the previous fairness certificate for
$P_0$ remains valid on $e$.
\end{proof}

A stochastic tester cannot in general distinguish an exact threshold from an
arbitrarily small violation with a uniform finite query bound.

\begin{proposition}[No finite uniform sharp-threshold test]
\label{arch:prop:zero-margin}
Fix $\tau\in(0,1/2)$ and $\delta\in(0,1/2)$.  On a graph containing one edge with
Bernoulli outputs, there is no tester with finite worst-case query bound that
uniformly distinguishes
\[
|p_u-p_v|\le\tau
\quad\text{from}\quad
|p_u-p_v|>\tau
\]
with success probability at least $1-\delta$ on every instance.
\end{proposition}

\begin{proof}
Take a fair instance on the boundary, e.g. $p_u=1/4$ and
$p_v=1/4+\tau$, and compare it with an unfair instance obtained by replacing
$p_v$ by $p_v+\eta$.  For any fixed finite query budget, the KL divergence between
the two possible transcript distributions tends to zero as $\eta\to0$.  Pinsker's
inequality therefore makes their total-variation distance arbitrarily small, so
no decision rule can maintain error at most $\delta$ in both worlds.
\end{proof}

Accordingly, the core problem uses a positive separation margin $\gamma>0$.


\section{Coupled Differential Execution}
\label{arch:sec:coupled}

A vertex-paired execution at vertex $v$ draws one random tape $R$ and evaluates both
versions:
\begin{equation}\label{eq:appendix-04-coupled-execution-001}
Y_{0,v}=P_0(v;R),\qquad Y_{1,v}=P_1(v;R),
\end{equation}
producing the raw pair $(Y_{0,v},Y_{1,v})$ and the derived delta
\begin{equation}\label{eq:appendix-04-coupled-execution-002}
Z_v=Y_{1,v}-Y_{0,v}.
\end{equation}
Thus one vertex-paired execution comprises two concrete program executions.
The statistical observation is theorem-specific: the Gaussian results use a
\emph{delta oracle} that reveals only $Z_v$, whereas the Bernoulli replay
results explicitly use the raw pair.  Edge, suite, and group replay reveal the
ordered raw outputs from all executions sharing the tape.  Fresh calls use
independent tapes.  Then $\E[Z_v]=\theta_v$.  Write
$\nu_v^2=\mathrm{Var}(Z_v)$.

\begin{proposition}[Differential variance]
\label{arch:prop:diffvar}
\begin{equation}\label{eq:appendix-04-coupled-execution-003}
\nu_v^2
=
\sigma_{1,v}^2+\sigma_{0,v}^2
-2\mathrm{Cov}(P_1(v;R),P_0(v;R)).
\end{equation}
If the two versions instead use independent random tapes, the variance of their
difference is $\widetilde\nu_v^2=\sigma_{1,v}^2+\sigma_{0,v}^2$.
\end{proposition}

\begin{proof}
Expand the variance of $Y_{1,v}-Y_{0,v}$.  Under independent replay the covariance
term vanishes.
\end{proof}

\begin{corollary}[Equal-variance correlation gain]
If both versions have variance $\sigma^2$ and correlation $\rho$ at every relevant
vertex, then
\begin{equation}\label{eq:appendix-04-coupled-execution-004}
\nu_v^2=2\sigma^2(1-\rho),
\qquad
\widetilde\nu_v^2=2\sigma^2.
\end{equation}
Thus every Gaussian delta-oracle bound homogeneous of degree one in the vertex
variances improves by factor $1-\rho$ under paired replay.
\end{corollary}

Common random numbers are not universally beneficial \cite{GlassermanYao1992}.
The Gaussian access model therefore exposes the actual paired-delta variance rather than
assuming positive correlation by fiat.


\section{What Version Information Actually Buys}
\label{arch:sec:consequences}

The Gaussian core lets us separate three kinds of version information that are
often conflated in regression testing: localization of where a change may occur,
statistical coupling of old and new executions, and numerical slack in the old
fairness certificate.  The first two can change query complexity directly.  The
third, by itself, does not improve the worst-case rate when patch effects are
unrestricted.

\subsection{A linear separation from patch localization}

Consider a fairness graph that is a matching of $M$ counterfactual edges and
assume unit paired-delta variance at every vertex.  Compare two testers with the
same paired execution oracle.  A \emph{patch-aware} tester receives a sound
impact set identifying the unique edge on which behavior may have changed.  A
\emph{location-unaware} tester is only promised that at most one edge changed,
but is not told which edge.

\begin{corollary}[Linear patch-locality separation]
\label{arch:cor:locality-separation}
Fix $\delta\in(0,1/2)$ and $\gamma>0$.  On the $M$-edge matching family above:
\begin{enumerate}[label=(\roman*)]
\item a patch-aware tester uses at most
\begin{equation}\label{eq:appendix-07-regression-consequences-001}
\frac{32\log(2/\delta)}{\gamma^2}+2
\end{equation}
paired-oracle calls; while
\item every adaptive location-unaware $\delta$-correct tester with finite expected
stopping time uses, on the all-neutral null instance,
\begin{equation}\label{eq:appendix-07-regression-consequences-002}
\E_0[T]\ge
\frac{8M\,\kl(1-\delta,\delta)}{\gamma^2}.
\end{equation}
\end{enumerate}
Consequently, for every fixed $\delta$, the worst-case separation is
$\Omega(M)$.
\end{corollary}

\begin{proof}
With a revealed sound impact edge, the relevant graph $H$ consists of one edge,
for which the one-edge closed form gives $\Psi(H)=4$.  Substitution into
\Cref{thm:upper} gives the first claim.

Without the edge location, correctness must hold against an alternative on every
one of the $M$ disjoint edges.  The relevant testing graph is therefore the full
matching, for which the matching closed form gives $\Psi=4M$.  Applying
\Cref{thm:lower} yields the second claim.  The ratio is linear in $M$ for fixed
confidence.
\end{proof}

The comparison does not give the patch-aware algorithm a stronger execution
oracle: its only additional information is the sound impact location.  Thus the
separation quantifies, in a clean family, the statistical value of software
change-impact information.

\subsection{Old certificate slack is not enough in the minimax sense}

Suppose an old edge $e=(u,v)$ is certified with gap
$g_{0,e}=\mu_{0,u}-\mu_{0,v}$ satisfying $|g_{0,e}|\le\tau_e$.  It is tempting to
expect large headroom $\tau_e-|g_{0,e}|$ to reduce the worst-case cost of testing
the new version.  This intuition is false without a restriction on the patch.

\begin{proposition}[No minimax gain from old slack alone]
\label{arch:prop:slack-no-free-lunch}
Consider a single fairness edge with Gaussian paired-delta observations of known
variances $\nu_u^2,\nu_v^2>0$.  The old gap $g_0$ and threshold $\tau$ may be known
exactly, with arbitrary slack $|g_0|\le\tau$.  A tester must distinguish
\begin{equation}\label{eq:appendix-07-regression-consequences-003}
\mathsf H_0:\ |g_0+\theta_u-\theta_v|\le\tau
\end{equation}
from
\begin{equation}\label{eq:appendix-07-regression-consequences-004}
\mathsf H_1:\ |g_0+\theta_u-\theta_v|\ge\tau+\gamma.
\end{equation}
For every adaptive $\delta$-correct tester with finite expected stopping time,
there is a boundary null instance for which
\begin{equation}\label{eq:appendix-07-regression-consequences-005}
\E[T]\ge
\frac{2\kl(1-\delta,\delta)}{\gamma^2}
\min_{q_u,q_v>0}
\left\{q_u+q_v:
\frac{\nu_u^2}{q_u}+\frac{\nu_v^2}{q_v}\le1
\right\}.
\end{equation}
In particular, the lower bound is independent of the old certificate slack.
\end{proposition}

\begin{proof}
Choose a sign $s\in\{-1,+1\}$ and a null patch effect whose edge difference is
\begin{equation}\label{eq:appendix-07-regression-consequences-006}
\theta_u^{(0)}-\theta_v^{(0)}=s\tau-g_0,
\end{equation}
so the new program lies exactly on the fairness boundary.  Construct an
alternative by increasing the signed edge difference by $\gamma$:
\begin{equation}\label{eq:appendix-07-regression-consequences-007}
\theta_u^{(1)}-\theta_v^{(1)}=s(\tau+\gamma)-g_0.
\end{equation}
The two hypotheses therefore differ by exactly $\gamma$ in the endpoint mean
difference, regardless of $g_0$.  Splitting this $\gamma$ shift adversarially
between the two endpoints and applying the same adaptive change-of-measure
calculation as in \Cref{thm:lower} gives the stated bound.
\end{proof}

\Cref{arch:prop:slack-no-free-lunch} does not say that historical certificates are
useless.  They can yield instance-dependent savings when the observed patch stays
well inside the boundary, and they become powerful under structural assumptions
that limit patch magnitude or direction.  The proposition says that such savings
cannot be claimed as a uniform minimax consequence of old slack alone.


\subsection{A two-axis information hierarchy}

The preceding Gaussian separation isolates the value of knowing \emph{where} a
patch may have changed behavior, while \Cref{thm:bernoulli-separation} isolates
the value of knowing \emph{how to compare} the two versions through replayable
randomness.  The two effects can be combined on a matching family.

Let \(G\) consist of \(M\) disjoint counterfactual edges
\(e_i=(u_i,v_i)\).  Fix \(p\in(0,1)\) and \(d\in(0,1-p)\).  The old version and
the null new version have Bernoulli parameter \(p\) at every vertex.  Alternative
\(i\) changes only \(u_i\) from \(p\) to \(p+d\).  Thus exactly one fairness edge
may regress, but its location is unknown.

\begin{corollary}[Unknown-location replay complexity]
\label{arch:cor:bernoulli-location}
Under common-uniform paired replay, a tester with no patch-location information
can be \(\delta\)-correct using
\begin{equation}\label{eq:appendix-07-regression-consequences-008}
M\left\lceil
\frac{\log\delta}{\log(1-d)}
\right\rceil
\end{equation}
paired-oracle calls.  Conversely, every adaptive \(\delta\)-correct tester with
finite expected stopping time satisfies on the all-neutral null
\begin{equation}\label{eq:appendix-07-regression-consequences-009}
\E_0[T]
\ge
\frac{M\,\kl(1-\delta,\delta)}
{(1-p)\log\!\left(\frac{1-p}{1-p-d}\right)}.
\end{equation}
Consequently, for fixed interior \(p\),
\begin{equation}\label{eq:appendix-07-regression-consequences-010}
Q_{\rm pair}^{\rm unknown}(M,d,\delta)
=
\Theta\!\left(
\frac{M\log(1/\delta)}{d}
\right)
\end{equation}
up to universal constant factors for, say, \(\delta\le 1/4\).
If a sound impact analysis identifies the affected edge in advance, the factor
\(M\) disappears by \Cref{thm:bernoulli-separation}.
\end{corollary}

\begin{proof}
For the upper bound, query each candidate vertex \(u_i\) for
\begin{equation}\label{eq:appendix-07-regression-consequences-011}
n=\left\lceil\frac{\log\delta}{\log(1-d)}\right\rceil
\end{equation}
paired replay rounds and reject if any old/new discordance is observed.  Under the
null, common-uniform replay makes every paired outcome identical, so the type-I
error is exactly zero even though \(M\) locations are inspected.  Under
alternative \(i\), only the samples at \(u_i\) are informative and each is
discordant with probability \(d\); hence the miss probability is
\((1-d)^n\le\delta\).

For the lower bound, fix alternative \(i\).  As in
\Cref{thm:bernoulli-separation}, a paired query to \(u_i\) contributes
\begin{equation}\label{eq:appendix-07-regression-consequences-012}
K_{\rm pair}
=
(1-p)\log\!\left(\frac{1-p}{1-p-d}\right)
\end{equation}
null-to-alternative KL, while every query to every other vertex has identical law
under the two hypotheses and contributes zero.  Adaptive transcript
change-of-measure therefore implies
\begin{equation}\label{eq:appendix-07-regression-consequences-013}
\E_0[N_{u_i}]K_{\rm pair}
\ge
\kl(1-\delta,\delta).
\end{equation}
The \(M\) candidate vertices are distinct.  Summing these inequalities over
\(i=1,\ldots,M\) yields the stated lower bound.
\end{proof}

\Cref{arch:cor:bernoulli-location} should not be read as a claim that searching for one
abnormal stochastic source is new; active sequential anomaly search is an
established information-theoretic problem \cite{CohenZhao2015}.  Its role here is
to expose a software-regression decomposition: sound change-impact information
removes the multiplicative location-search cost, while replayable old/new
randomness changes the small-patch dependence from quadratic to linear.


\subsection{When old certificate headroom does become useful}

\Cref{arch:prop:slack-no-free-lunch} rules out a minimax gain from old fairness slack
when a patch is unrestricted.  The conclusion changes if program analysis
provides a \emph{sound patch-effect envelope}.  For edge \(e=(u,v)\), let
\begin{equation}\label{eq:appendix-07-regression-consequences-014}
b_e
\ge
|(\mu_{1,u}-\mu_{0,u})-(\mu_{1,v}-\mu_{0,v})|
=
|\theta_u-\theta_v|.
\end{equation}
Such an envelope is deterministic side information about the patch; the theorem
does not prescribe how it is obtained.

\begin{proposition}[Zero-query certificate carryover]
\label{arch:prop:certificate-pruning}
Let the old edge gap \(g_{0,e}\) and fairness threshold \(\tau_e\) be known.  If
\begin{equation}\label{eq:appendix-07-regression-consequences-015}
|g_{0,e}|+b_e\le\tau_e,
\end{equation}
then edge \(e\) is fair under \(P_1\) and requires no stochastic execution to
re-certify.
\end{proposition}

\begin{proof}
Using \(g_{1,e}=g_{0,e}+\theta_u-\theta_v\) and the triangle inequality,
\begin{equation}\label{eq:appendix-07-regression-consequences-016}
|g_{1,e}|
\le
|g_{0,e}|+|\theta_u-\theta_v|
\le
|g_{0,e}|+b_e
\le
\tau_e.
\end{equation}
\end{proof}

Define the \emph{residual regression graph} \(H^\star\) to contain exactly the
edges for which
\begin{equation}\label{eq:appendix-07-regression-consequences-017}
|g_{0,e}|+b_e>\tau_e.
\end{equation}
All other fairness comparisons are discharged by
\Cref{arch:prop:certificate-pruning} before any paired execution is performed.

\begin{corollary}[Envelope-pruned Gaussian regression testing]
\label{arch:cor:envelope-gaussian}
Suppose paired deltas are Gaussian with known vertex variances \(\nu_v^2\), the
old gaps \(g_{0,e}\) are known, and the tester faces the separated promise
\begin{equation}\label{eq:appendix-07-regression-consequences-018}
\mathsf H_0:\ |g_{1,e}|\le\tau_e\quad\forall e,
\end{equation}
versus
\begin{equation}\label{eq:appendix-07-regression-consequences-019}
\mathsf H_1:\ |g_{1,e}|\ge\tau_e+\gamma
\quad\text{for some }e,
\end{equation}
for \(\gamma>0\).  If \(m^\star=|E(H^\star)|>0\), a non-adaptive tester exists
using at most
\begin{equation}\label{eq:appendix-07-regression-consequences-020}
\frac{8\log(2m^\star/\delta)}{\gamma^2}
\Psi(H^\star,\nu)
+
|V(H^\star)|
\end{equation}
paired-oracle calls with error probability at most \(\delta\).  If
\(m^\star=0\), zero paired-oracle calls suffice.
\end{corollary}

\begin{proof}
Edges outside \(H^\star\) are fair by
\Cref{arch:prop:certificate-pruning}.  On the residual graph, allocate paired samples
exactly as in \Cref{thm:upper}.  For each residual edge \(e=(u,v)\), estimate
\(\theta_u-\theta_v\), add the known offset \(g_{0,e}\), and reject if
\begin{equation}\label{eq:appendix-07-regression-consequences-021}
|g_{0,e}+\widehat\theta_u-\widehat\theta_v|
\ge
\tau_e+\gamma/2.
\end{equation}
The allocation makes every residual edge-difference estimation error at most
\(\gamma/2\) simultaneously except with probability \(\delta\).  On this event,
all fair instances are accepted and every alternative separated by \(\gamma\) is
rejected.  The query count is the Gaussian allocation bound applied only to
\(H^\star\).
\end{proof}

This result gives a precise role to retained certificate headroom.  Slack is not
useful merely because it exists; it becomes computationally useful when coupled
to sound information constraining what the patch can do.  Better static bounds
shrink \(H^\star\), which can reduce both the number of tested comparisons and
the graph allocation complexity \(\Psi\).


\section{From the Core Problem to Full Fairness Regression}
\label{arch:sec:open}

The preceding theorems deliberately isolate the cleanest core.  Several extensions
are required for the final paper.

\paragraph{Retained fairness headroom.}
For edge $e$, the old version has a certified gap $g_{0,e}$ and threshold $\tau_e$.
The quantity $\tau_e-|g_{0,e}|$ is reusable statistical headroom.  A full tester
should spend little effort on an edge whose old certificate is far from the
boundary and whose observed patch delta points inward.

\paragraph{Unknown delta variance.}
The oracle quantity $\Psi(H,\nu)$ assumes known $\nu_v^2$.  A practical sequential
tester must estimate variance and maintain anytime-valid confidence while adapting
its sample allocation.

\paragraph{Imperfect change impact.}
A sound but imprecise impact set wastes tests on false-positive vertices.  A
separate controlled-unsoundness model can reserve a small sentinel budget for
behavior outside the reported impact set.

\paragraph{Multi-patch amortization.}
For a release sequence $P_0\to P_1\to\cdots\to P_T$, fairness certificates and
paired-delta evidence should be reusable across several versions.  The objective
is an amortized bound governed by cumulative impacted structure rather than
$T|E|$.

\paragraph{Partial synchronized replay.}
Between edgewise four-way replay and full synchronized-suite replay lies an
intermediate oracle in which only declared groups of fairness-relevant inputs can
share a random tape.  Let $\mathcal G=\{S_1,\ldots,S_k\}$ be the legal replay
groups, each of size at least two, with query costs $c_i>0$.  Two incidence
matrices arise naturally.  The edge-by-group witness matrix
\[
A_{e,i}=\mathbf{1}\{e\subseteq S_i\}
\]
governs certified four-way witnesses, whereas the active-vertex matrix
\[
B_{x,i}=\mathbf{1}\{x\in S_i\},
\qquad x\in V_+,
\]
governs raw transcript information in the homogeneous Bernoulli family.

For the witness construction, define
\[
C_{\rm edge}
=
\min_{a\ge0}
\left\{
c^\top a: Aa\ge\mathbf{1}
\right\}.
\]
If
\[
L=
\left\lceil
\frac{\log\delta}{\log(1-d)}
\right\rceil,
\qquad
n_i=\lceil L a_i^\star\rceil,
\]
then querying group $S_i$ exactly $n_i$ times gives every fairness edge at least
$L$ independent legal witness opportunities.  The resulting fixed design has
zero type-I error in the frozen common-patch model, pointwise type-II error at
most $\delta$, and declared cost at most
\[
L C_{\rm edge}
+
\sum_{i:a_i^\star>0} c_i.
\]

The raw-output lower geometry is different.  For an oriented alternative whose
changed endpoint is $x$, every legal group containing $x$ has one-query
null-to-alternative KL
\[
K_{\rm sync}
=
r\log\!\left(\frac r{r-d}\right),
\qquad r=1-p-s,
\]
regardless of whether the second synchronized vertex is the fairness partner.
Thus transcript change of measure reduces to the active-vertex fractional cover
\[
C_{\rm vert}
=
\min_{q\ge0}
\left\{
c^\top q:Bq\ge\mathbf{1}
\right\},
\]
yielding the necessary expected-cost floor
\[
\frac{\kl(1-\delta,\delta)}{K_{\rm sync}}C_{\rm vert}.
\]
Every feasible edge-cover vector is also feasible for the active-vertex cover,
so $C_{\rm vert}\le C_{\rm edge}$.

\begin{proposition}[Separation of the P17 LP geometries]
\label{arch:prop:p17-geometry-separation}
The two fractional-cover geometries need not be comparable within any finite
universal factor.  More precisely: (i) if every fairness edge has endpoints
with identical legal-group incidence signatures, then
$C_{\rm vert}=C_{\rm edge}$; (ii) for $H=K_n$ with the legal groups equal to
its edges and unit group costs,
$C_{\rm vert}=n/2$ and $C_{\rm edge}=\binom n2$, so their ratio is $n-1$;
and (iii) there is a legal partial-replay instance with
$C_{\rm vert}<\infty$ but $C_{\rm edge}=\infty$.
\end{proposition}

\begin{proof}
For (i), identical endpoint incidence makes the edge-witness row
$\mathbf 1\{u,v\in S_i\}_i$ identical to either endpoint's active-vertex
row; equality propagates over each connected fairness component, and duplicate
constraints may be removed.  For (ii), pair groups cover only their own edge, so
every edge variable in the witness LP is at least one.  Summing the $n$
active-vertex constraints counts every pair variable twice and gives
$C_{\rm vert}\ge n/2$, while assigning weight $1/(n-1)$ to every pair
attains equality.  For (iii), take the path $0$--$1$--$2$ with legal groups
$\{0,2\}$ and $\{1,2\}$ at unit cost.  The active-vertex LP has value $2$,
but edge $\{0,1\}$ belongs to no legal group, so the witness LP is infeasible.
\end{proof}

Proposition~\ref{arch:prop:p17-geometry-separation} concerns only the two LP
geometries; it is not by itself a testing-complexity separation.  Raw outputs,
however, permit a stronger construction in the same homogeneous family.

\begin{proposition}[Active-vertex raw-contrast schedule]
\label{arch:prop:p17-active-contrast}
Assume every vertex appearing in a legal replay group has old threshold $p$
and, under the null, new threshold $p+s$.  Under each alternative, exactly one
fairness-active endpoint $x$ instead has new threshold $p+s+d$, with
$p,s,d>0$ and $p+s+d<1$.  Suppose $C_{\rm vert}<\infty$.  Let $q^\star$
solve the active-vertex LP, set
\[
L=\left\lceil\frac{\log\delta}{\log(1-d)}\right\rceil,
\qquad
n_i=\lceil Lq_i^\star\rceil,
\]
and query $S_i$ exactly $n_i$ times with independent tapes across queries.
Within a query, define $D_v=P_1(v;U)-P_0(v;U)$ and reject if the vector
$(D_v:v\in S_i)$ is nonconstant.  The test has zero type-I error and
pointwise type-II error at most $\delta$, with declared cost at most
\[
L C_{\rm vert}+\sum_{i:q_i^\star>0}c_i.
\]
Consequently the corrected information lower bound and this executable upper
construction use the same active-vertex combinatorial geometry in the frozen
homogeneous single-changed-endpoint family.
\end{proposition}

\begin{proof}
Under the null, common-uniform replay gives
$D_v=\mathbf 1\{U\in(p,p+s]\}$ for every member of every legal group, so
the within-group delta vector is constant pathwise and rejection is impossible.
Under an alternative changing $x$, any legal group containing $x$ contains at
least one distinct member $y$ because legal groups have size at least two.
The alternative changes only $x$, hence
\[
D_x\ne D_y
\quad\Longleftrightarrow\quad
U\in(p+s,p+s+d],
\]
which has probability $d$.  For every active $x$, feasibility of $q^\star$
and upward rounding give
\[
\sum_{i:x\in S_i}n_i
\ge L\sum_{i:x\in S_i}q_i^\star
\ge L.
\]
Independent tapes across group queries therefore make the miss probability for
alternative $x$ at most $(1-d)^L\le\delta$.  Finally,
$\lceil Lq_i^\star\rceil\le Lq_i^\star+1$ on the positive support,
which gives the stated cost bound.
\end{proof}

Let $Q^\star_{\rm hom}$ denote the infimum, over tests with type-I and
pointwise type-II error at most $\delta$ on this homogeneous family, of the
worst-case expected declared query cost.  The transcript lower bound applies in
particular to the null expected cost, while the construction above has a
deterministic declared cost.  Combining the two therefore yields
\[
\frac{\kl(1-\delta,\delta)}{K_{\rm sync}}C_{\rm vert}
\;\le\;
Q^\star_{\rm hom}
\;\le\;
L C_{\rm vert}+\sum_{i:q_i^\star>0}c_i.
\]
This identifies the active-vertex geometry for the frozen homogeneous family
up to finite-sample constants and rounding.  It also explains why the
path/reference construction of Proposition~\ref{arch:prop:p17-geometry-separation}
remains testable even though its four-way witness LP is infeasible.

The scope is essential.  For a heterogeneous fairness-preserving null, patch
deltas at two vertices can differ pathwise without constituting a fairness
regression, so the zero-type-I contrast argument need not hold.

\begin{remark}[A scope counterexample]
\label{arch:rem:p17-heterogeneous-scope}
Take a single fairness edge $\{0,1\}$ with tolerance $\tau=0.20$. Let the
old Bernoulli thresholds be $(0.20,0.40)$ and the new thresholds be
$(0.30,0.45)$. The new-version gap is $0.15\le\tau$, so the fairness null is
satisfied. Yet at shared tape $U=0.25$ the old outputs are $(0,1)$ and the new
outputs are $(1,1)$, giving patch deltas $(1,0)$. The homogeneous raw-contrast
predicate therefore rejects a fairness-preserving heterogeneous null.
In fact, the rejection event is not a measure-zero witness: the patch deltas
are unequal for
\[
U\in(0.20,0.30]\cup(0.40,0.45],
\]
whose exact probability under the common-uniform tape is $0.10+0.05=0.15$.
This is an analytic property of the frozen specification, not an empirical
false-positive estimate. Consequently Proposition~\ref{arch:prop:p17-active-contrast} cannot be extended to
the general composite null without additional controls or a different decision
statistic.
\end{remark}

The fully heterogeneous composite-null setting remains open, but it is no
longer entirely uncharacterized.  \Cref{thm:control-replay} gives a positive
two-class bipartite extension: when legal replay groups contain
matched-background controls, first-order information and a fractional
$C_{\rm ctrl}$ geometry reappear; complete control coverage restores the linear
minimax order, while without controls the fixed-tolerance boundary channel is
quadratic.  In mixed topologies a boundary/control hybrid matches the
small-effect orders for each fixed finite topology.  A fixed nonadaptive
boundary allocation provably pays $|U|\log |U|/d^2$, while an unbounded-horizon
sequential likelihood-ratio test removes the multiplicative $\log |U|$ from
expected null cost, and Markov truncation retains $\Theta(|U|/d^2)$ topology
scaling under a deterministic adaptive cost cap at fixed confidence.  Exact
adaptive constants and optimal confidence dependence remain open.  The broader
open problem is to move
beyond known control classes to multiple or unknown background levels and more
general fairness graphs.  Common random
numbers and their use in multiple-system
comparisons are established simulation tools
\cite{GlassermanYao1992,YangNelson1991}; the claim here is the resulting
version-aware fairness-regression query geometry, not the invention of common
randomness.

\paragraph{Composite-null action information (P20, open).}
A central unresolved difficulty is that a fairness-preserving nuisance parameter
$\eta$ belongs to the \emph{whole} null program, not to each replay action
independently. For a frozen family of replay actions $S$ and oriented
alternatives $a$, define the exact null-to-alternative information
\[
I_{a,S}(\eta)
=
D_{\mathrm{KL}}\!\left(
P_{0,\eta}^{S}\,\middle\|\,P_{a,\eta}^{S}
\right).
\]
Under the usual stopped-transcript change-of-measure assumptions, every
level-$\delta$ policy obeys, separately for every fixed $\eta$ and $a$,
\[
\sum_S
\mathbb{E}_{0,\eta}[N_S]\,I_{a,S}(\eta)
\ge
\kl(1-\delta,\delta).
\]
Consequently one obtains a necessary uniform-null cost floor by maximizing
the \emph{pointwise-nuisance} information-cover minimum over $\eta$.
This must not be confused with minimizing cost for one fixed allocation
subject to every nuisance constraint simultaneously: expected adaptive
action counts may depend on the true nuisance.

Even for two actions and two nuisance values, the distinction is strict.
With unit costs, information rows $(1,\varepsilon)$ and
$(\varepsilon,1)$, and unit required information, both pointwise
minima equal $1$, whereas the shared fixed-information design costs
$2/(1+\varepsilon)$. This is an abstract LP separation, not yet a
realizable software-oracle or optimal-testing separation.
The detailed quantifier audit and finite fixtures are maintained in the
replication repository; P20 remains open. Composite controlled sensing,
restricted-action error exponents, and active testing algorithms are already
studied in the archival literature
\cite{DeshmukhVeeravalliBhashyam2021,HsuWang2025,GanJiaLiTayur2026}.
Our remaining question is which such information patterns can arise from
\emph{legal same-tape executions of evolving software} and which admit
matching sound regression-certification procedures.

\paragraph{Adaptivity and the Gaussian logarithmic gap.}
The Gaussian upper bound in \Cref{thm:upper} pays
$\log(|E(H)|/\delta)$ for simultaneous edge control, while the current adaptive
null lower bound in \Cref{thm:lower} contains only $\log(1/\delta)$.  This gap
should not be framed as a generic discovery that sequential sampling can remove
dimension penalties: sequential sparse-recovery work already establishes large
adaptive savings relative to nonsequential coordinate-wise procedures
\cite{MalloyNowak2014}, and property testing also exhibits strict adaptivity
hierarchies \cite{CanonneGur2018,Hurwitz2011}.  The unresolved question here is
more specific: whether shared vertex observations and fairness-graph incidence
permit an adaptive TFOESUU tester whose cost is governed by $\Psi(H,\nu)$
without the full edgewise union-bound factor, or whether some graph family forces
an additional structural logarithm.  Any future claim must separate expected
null cost, alternative cost, replay rounds, and concrete program executions.


\section{Computational Validation and Experimental Protocol}
\label{arch:sec:evaluation}

The empirical component of this paper is intentionally a \emph{theorem
falsification and reproducibility layer}, not a leaderboard study. Its purpose
is to check whether executable implementations reproduce the exact structural
predictions of proved results, to search for counterexamples outside proved
scope, and to measure finite-instance constants only after the corresponding
theorem gate is frozen. We therefore distinguish deterministic specification
results from stochastic publication experiments throughout.

\subsection{Research questions}

\begin{description}[leftmargin=1.8cm,labelwidth=1.45cm,style=nextline]
\item[\textbf{RQ1}] \textbf{Graph allocation.}
Do executable Gaussian allocation instances reproduce the predicted dependence
on the uncertainty-weighted graph quantity $\Psi(H,\nu)$ across graph families,
variance profiles, and effect sizes?

\item[\textbf{RQ2}] \textbf{Replay advantage.}
Under matched old/new, four-way, and suite replay access, do finite-instance
runs reproduce the theorem-predicted small-effect laws, including the P18
transition from $d^{-1}$ behavior at $\tau=O(d)$ to $d^{-2}$ behavior at fixed
positive fairness tolerance?

\item[\textbf{RQ3}] \textbf{Partial synchronized replay.}
Does executable group scheduling reproduce the $C_{\rm vert}$ geometry in the
homogeneous family, the new $C_{\rm ctrl}$ geometry when matched-background
controls are available, and the predicted quadratic fallback when they are not?

\item[\textbf{RQ4}] \textbf{Reproducibility and robustness.}
Are theorem-facing conclusions invariant to regenerated manifests, exact seeds,
graph materialization, implementation round-trips, and clean-checkout execution?
\end{description}

\subsection{Frozen specification datasets}

\Cref{arch:tab:benchmark-inventory} records the complete checked-in specification
inventory. These are synthetic \emph{inputs with exact ground truth}, not
observational datasets and not measured outcomes. They are deliberately compact
so reviewers can regenerate every configuration and audit theorem/code
correspondence.

\begin{table}[H]
\centering
\caption{Frozen benchmark/specification inventory. Row counts describe
deterministic test configurations, not empirical sample sizes.}
\label{arch:tab:benchmark-inventory}
\normalsize
\begin{tabular}{@{}p{0.34\textwidth}r p{0.48\textwidth}@{}}
\toprule
Artifact & Rows & Primary role \\
\midrule
Gaussian allocation grid & 384 & Graph families, variance profiles, effect-size and allocation stress. \\
Paired replay grid & 54 & Paired-versus-independent Bernoulli replay specifications. \\
Executable software subjects & 32 & Randomized software configurations exercising the execution interface. \\
Four-way counterfactual replay & 144 & Common-effect and differential-effect regimes for P15. \\
P18 threshold phase grid & 24 & Zero, matched-scale, and fixed positive fairness tolerances across four differential margins. \\
P19 control-replay fixtures & 6 & Exact control-cover feasibility, weighted-cover values, and boundary-only quadratic structure. \\
P17 partial replay stress grid & 168 & Group size, graph family, group cost, and small-effect stress. \\
P17 structural fixtures & 8 & Exact equality, finite-gap, and infeasible-witness constructions. \\
P17 scope counterexample & 1 & Heterogeneous fair null that falsifies naive raw-contrast generalization. \\
\midrule
\textbf{Total} & \textbf{821} & Deterministic, regenerable specification rows. \\
\bottomrule
\end{tabular}
\end{table}

\subsection{Public evaluation populations}

The theorem-specification manifests are complemented by two noninterchangeable
seven-dataset protocols: a theorem-aligned simulation and a direct executable
software-version comparison. These datasets do not establish the theorems.
They supply realistic covariate populations, and the second protocol also uses
aligned observed protected attributes. The raw datasets are never vendored
into the repository. Their source metadata are frozen in
\texttt{data/public\_benchmark\_registry\_v1.jsonl}, and publication runs must record a
local content hash before any result is admitted to the paper.

\begin{table}[H]
\centering
\caption{Planned public benchmark populations. The executable old/new software
versions are constructed locally from each dataset under the frozen protocol;
no pretrained model or external model API is used.}
\label{arch:tab:public-datasets}
\normalsize
\begin{tabular}{@{}p{0.25\textwidth}p{0.25\textwidth}p{0.20\textwidth}p{0.20\textwidth}@{}}
\toprule
Dataset & Task & Protected attributes & Provenance \\
\midrule
Adult Census Income & Binary classification & sex, race & UCI \cite{AdultDatasetUCI} \\
Bank Marketing & Binary classification & age & UCI \cite{BankMarketingUCI} \\
Statlog German Credit & Binary classification & sex, age & UCI \cite{GermanCreditUCI} \\
ProPublica COMPAS & Binary classification & race, sex & ProPublica \cite{ProPublicaCOMPAS2016} \\
Law School GPA & Regression to thresholded decision & race, gender & AIF360 \cite{BellamyEtAl2019} \\
MEPS Panel 21 & Binary classification & race & AIF360/AHRQ \cite{BellamyEtAl2019} \\
Default of Credit Card Clients & Binary classification & sex, age & UCI \cite{DefaultCreditUCI} \\
\bottomrule
\end{tabular}
\end{table}

\paragraph{Theorem-aligned simulation protocol.}
The simulation protocol fits a simple local baseline predictor on non-protected
features, converts its score to an interior Bernoulli probability, and defines
the old program by a common-uniform randomized decision. In this evaluator,
each evaluation probability is duplicated into a controlled two-endpoint edge;
the prepared artifact does not store or mutate protected-attribute values and
therefore does not materialize an empirical counterfactual pair. The new program
is produced by a frozen patch specification with a common shift and, when
required by the experiment, a differential endpoint shift. Thus public
covariates determine a heterogeneous probability population, while the graph
and counterfactual endpoints remain a synthetic theorem-aligned wrapper. This
construction is not evidence that the source dataset supplies valid causal
counterfactual worlds or satisfies the paper's fairness assumptions.

\paragraph{Historical executable-software protocol.}
A separately versioned observational-context artifact preserves aligned
training and evaluation rows, protected columns, labels, source-row positions,
and non-protected score-model inputs. A score companion records the exact
prefit training and evaluation probabilities without reconstructing discarded
dataframe types. The executable subject is a pinned direct parent/child pair of
Fairlearn's threshold optimizer: the child repairs multi-column protected-group
handling while leaving the prediction and randomness implementations fixed.
For each dataset and revision, the frozen primary estimand is the row-wise
difference between exact randomized-prediction PMFs obtained from the native
protected columns and from an injective encoding of the same ordered tuple.
Single-attribute datasets are negative controls. This measures a software
behavioral regression under observed groups; it is neither a causal fairness
effect nor a population-level fairness estimate. The immutable input bridge,
support rules, complete seven-dataset census, and source/runtime gates were
frozen before execution.

The recorded experiment passed the specified provenance checks and contains the complete fourteen-record
old/new census.  Table~\ref{tab:historical-census} reports every
preregistered primary estimand.  The parent revision disagreed with its injective
canonical control on Adult Census Income and Default of Credit Card Clients,
but produced exact zeros on COMPAS, German Credit, and Law School GPA.  Both
single-attribute controls were also exact zeros.  The repaired revision matched
the canonical control exactly on all seven datasets.  These are deterministic
software-behavior results on the frozen evaluation rows; they do not estimate a
causal fairness effect or population generalization.

\begin{table}[H]
\centering
\caption{Historical protected-group encoding comparison across seven input populations. Entries are maxima and means of exact row-wise absolute differences between native and canonical randomized-prediction probability masses; zero indicates agreement on the evaluated rows.}
\label{tab:historical-census}
\small
\begin{tabular}{@{}lrrrr@{}}
\toprule
& \multicolumn{2}{c}{Parent $|\Delta\mathrm{PMF}|$} & \multicolumn{2}{c}{Repaired $|\Delta\mathrm{PMF}|$} \\
\cmidrule(lr){2-3}\cmidrule(l){4-5}
Dataset & Max & Mean & Max & Mean \\
\midrule
Adult Census Income & 1.0000 & 0.0552 & 0.0000 & 0.0000 \\
Bank Marketing & 0.0000 & 0.0000 & 0.0000 & 0.0000 \\
Statlog German Credit & 0.0000 & 0.0000 & 0.0000 & 0.0000 \\
ProPublica COMPAS & 0.0000 & 0.0000 & 0.0000 & 0.0000 \\
Law School GPA & 0.0000 & 0.0000 & 0.0000 & 0.0000 \\
MEPS Panel 21 & 0.0000 & 0.0000 & 0.0000 & 0.0000 \\
Default of Credit Card Clients & 1.0000 & 0.1023 & 0.0000 & 0.0000 \\
\bottomrule
\end{tabular}
\vspace{2pt}
\begin{minipage}{0.98\linewidth}\footnotesize
The primary analysis uses deterministic exact PMFs. No confidence interval applies; four preselected common-tape seeds served as implementation-correspondence checks. All datasets, including exact zeros, are reported.
\end{minipage}
\end{table}


Prepared inputs and experiment code carry separate provenance.  A later clean
execution commit may consume an earlier prepared artifact only if a fail-closed
compatibility check verifies identical Git objects for the preparation script,
frozen public registry, and public dependency contract.  Every admitted result
then records both commits and the digest of that compatibility contract; any
preparation-surface change requires re-preparation rather than metadata edits.

\paragraph{Simulation unit and execution accounting.}
The public GPU evaluator implements a controlled \emph{Bernoulli threshold-law
simulation}, not a wall-clock profile of an instrumented operating system.
Independent and old--new vertex-paired baselines sample their exact binomial
sufficient statistics. In the frozen single-changed-edge experiment, the
four-way, suite, and localized tester share the same common-tape threshold
witness trace; all other edges are null and cannot trigger the pathwise witness.
Their reported \emph{concrete program executions} therefore count the calls
required by the declared replay schedule, rather than GPU kernel executions
actually performed by this sufficient-statistic simulator. Shared witness
errors are deliberately correlated across these three access models and
must not be interpreted as three independent empirical replicates. No
generalization to arbitrary stochastic software, uncontrolled nondeterminism,
or heterogeneous patch nulls follows from this experiment.
\subsection{Experimental settings and baselines}

Every stochastic publication run is keyed by the repository commit, a frozen
JSONL row, an explicit seed, the execution backend, and a machine-environment
record. CPU experiments are used for exact or moderate-size checks; GPU sweeps
are reserved for high-trial finite-instance studies after the relevant theorem
and manifest are frozen. No pretrained model or external API is used.

\paragraph{Two-workstation execution environment (preflight).}
The planned local study uses two independent workstations, each independently
identified by the operating system as containing one NVIDIA GeForce RTX 5070 Ti
GPU with 16\,GiB-class memory (16{,}303\,MiB reported per device).
Both hosts reported NVIDIA driver 591.86 and driver-advertised CUDA 13.1
compatibility on October 8, 2026. The first host uses Windows with Python
3.11 available; the second runs Ubuntu 22.04.5 under WSL2 with Python 3.12.2
available. Hardware inventory alone is not outcome evidence. Engineering
smokes and preflights are excluded from submission-facing tables unless their
complete manifests pass the corresponding release gate. The admitted P12/P15
manifest binds source, runners, specifications, trial budgets, and result
hashes, but does not contain a complete PyTorch/CUDA environment fingerprint;
we therefore make no environment-sensitive performance claim. The pinned executable-version protocol instead uses a
CPU-only Python~3.8.20 historical runtime; GPU acceleration is scientifically
irrelevant to that subject.

\paragraph{Independent workers and reproducibility controls.}
The two GPUs are separate devices on different machines, \emph{not} one
shared-memory two-GPU accelerator. A Windows controller on the first host
can dispatch an SSH command to the second host, whose Ubuntu WSL2 checkout
performs computations locally; the second host's Windows D: volume is
designated for persistent datasets, logs, and result artifacts. Before a
shard begins, both workers must have the same clean Git commit and
a checked, frozen manifest. Shards are assigned disjoint configuration
identifiers with explicit seeds, device information, input hashes, and
trial budgets. Results may be merged only after verifying those keys and
recording failures as well as successes. Trial budgets are staged from
$10^3$ exploratory replications through $10^4$ and $10^5$, with larger
rare-error sweeps contingent on passing earlier validation gates.
Per-worker GPU runtime and memory are not equated with the
\emph{concrete program execution} complexity used in theorems.
No two-host performance or publication-scale empirical result is asserted;
the admitted 100k campaign is used only for exact oracle-law validation.

The comparison set is access-model based rather than model based:
\begin{enumerate}[leftmargin=1.6em,itemsep=2pt]
\item fresh/independent execution where no old--new coupling is available;
\item old/new vertex pairing under a common tape;
\item four-way edge replay when both endpoints can share that tape;
\item synchronized-suite replay when one tape covers all active inputs; and
\item declared-group partial replay, including the $C_{\rm edge}$ witness
schedule and the $C_{\rm vert}$ raw-contrast schedule where its assumptions hold.
\end{enumerate}
Primary quantities are concrete program executions, replay rounds, empirical
type-I/type-II error for stochastic runs, and normalized scheduling cost.
Whenever a theorem is asymptotic in $d$, the finite-instance study varies $d$
over the frozen manifest rather than fitting an exponent to hand-selected
points.

\subsection{Current deterministic audit results}

Before any publication-scale stochastic sweep, the repository already supports
exact deterministic checks. \Cref{arch:tab:deterministic-evidence} reports only
facts that follow from the committed specifications and executable theorem
checks; it intentionally does not report fabricated performance numbers.

\begin{table}[H]
\centering
\caption{Current deterministic evidence available before stochastic publication
runs. Pass condition states what the executable audit must verify.}
\label{arch:tab:deterministic-evidence}
\normalsize
\begin{tabular}{@{}p{0.23\textwidth}p{0.33\textwidth}p{0.35\textwidth}@{}}
\toprule
Audit target & Frozen evidence & Pass condition \\
\midrule
P17 stress geometry & 168 parameter rows & Materialized group/cost geometry is valid; $C_{\rm vert}\le C_{\rm edge}$; rounded active-vertex exposure attains the required miss bound. \\
P17 structural gap & 8 exact fixtures & One equality case, six clique cases with $C_{\rm edge}/C_{\rm vert}=n-1$, and one finite-$C_{\rm vert}$/infeasible-$C_{\rm edge}$ case agree with the LP implementations. \\
P17 scope boundary & 1 heterogeneous fair-null fixture & New-version gap is $0.15\le\tau=0.20$, while the homogeneous raw-contrast predicate rejects on tape set $(0.20,0.30]\cup(0.40,0.45]$ of exact measure $0.15$. \\
P18 tolerance phase & 24 parameter rows & Every row yields positive exact transcript KL and a valid theorem schedule; zero-tolerance rows recover the linear replay regime while fixed positive tolerances probe the quadratic regime. \\
P19 control geometry & 6 exact fixtures & Fractional control-cover feasibility and expected costs match the frozen specifications; the boundary-only fixture remains information-observable but loses the linear control channel. \\
Manifest integrity & 821 total rows & All frozen artifacts exist, IDs are unique, row counts meet the frozen inventory, and regenerators round-trip exactly. \\
\bottomrule
\end{tabular}
\end{table}

\begin{figure}[H]
\centering
\begin{tikzpicture}[x=0.855cm,y=0.68cm,
  axis/.style={draw=black!55,line width=0.65pt},
  guide/.style={draw=black!15,line width=0.45pt},
  focalLine/.style={draw=focalteal,line width=1.35pt},
  riskLine/.style={draw=risktone,line width=3.2pt,line cap=round},
  pt/.style={circle,fill=focalteal,inner sep=1.8pt},
  lab/.style={font=\scriptsize,text=black!78},
  small/.style={font=\scriptsize,text=black!65}
]

% Panel (a): exact structural ratio.
\begin{scope}[shift={(0,0)}]
\node[lab,anchor=west,font=\small\bfseries,fill=white,inner sep=1pt] at (0.85,6.58) {(a) Exact clique gap};
\draw[axis,->] (0.5,0.75) -- (7.7,0.75);
\draw[axis,->] (0.5,0.75) -- (0.5,6.65);
\foreach \x/\n in {1.4/3,2.4/4,3.4/5,4.4/6,5.4/7,6.4/8}{
  \draw[axis] (\x,0.68)--(\x,0.82);
  \node[small,anchor=north] at (\x,0.58) {\n};
}
\foreach \y/\v in {1.5/1,2.3/2,3.1/3,3.9/4,4.7/5,5.5/6,6.3/7}{
  \draw[guide] (0.5,\y)--(7.35,\y);
  \node[small,anchor=east] at (0.36,\y) {\v};
}
\draw[focalLine] (1.4,2.3)--(2.4,3.1)--(3.4,3.9)--(4.4,4.7)--(5.4,5.5)--(6.4,6.3);
\foreach \x/\y in {1.4/2.3,2.4/3.1,3.4/3.9,4.4/4.7,5.4/5.5,6.4/6.3}{\node[pt] at (\x,\y) {};}
\node[lab,anchor=north] at (4.1,0.08) {Clique size $n$};
\node[lab,rotate=90,anchor=south] at (-0.32,3.45) {$C_{\rm edge}/C_{\rm vert}$};
\node[small,anchor=west,text=focalteal] at (4.7,2.65) {exactly $n-1$};
\end{scope}

% Panel (b): exact scope obstruction.
\begin{scope}[shift={(9.0,0)}]
\node[lab,anchor=west,font=\small\bfseries,fill=white,inner sep=1pt] at (0,6.58) {(b) Heterogeneous null breaks raw contrast};
\draw[axis,->] (0.45,3.4)--(7.9,3.4);
\foreach \x in {0.45,1.90,2.63,3.35,3.72,7.70}{
  \draw[axis] (\x,3.31)--(\x,3.49);
}
\node[small,anchor=north] at (0.45,3.18) {0};
\node[small,anchor=north] at (1.90,3.18) {0.20};
\node[small,anchor=north] at (2.63,2.91) {0.30};
\node[small,anchor=north] at (3.35,3.18) {0.40};
\node[small,anchor=north] at (3.72,2.91) {0.45};
\node[small,anchor=north] at (7.70,3.18) {1};
\draw[riskLine] (1.90,3.4)--(2.63,3.4);
\draw[riskLine] (3.35,3.4)--(3.72,3.4);
\node[small,text=risktone,anchor=south] at (2.27,3.67) {reject};
\node[small,text=risktone,anchor=south] at (3.54,3.67) {reject};
\node[lab,anchor=north] at (4.0,2.40) {shared tape $U$};
\node[lab,anchor=west] at (0.45,1.72) {Fairness null: $|0.30-0.45|=0.15\le\tau=0.20$};
\node[lab,anchor=west,text=risktone] at (0.45,1.15) {Raw-contrast rejection measure: $0.10+0.05=0.15$};
\node[small,anchor=west,text width=6.2cm,align=left] at (0.45,0.52)
  {The homogeneous $C_{\rm vert}$ test is invalid for the unrestricted composite null.};
\end{scope}
\end{tikzpicture}
\caption{Exact deterministic P17 audit evidence.  (a) For clique $K_n$ with legal
replay groups equal to edges and unit costs, the witness geometry satisfies
$C_{\rm edge}/C_{\rm vert}=n-1$, giving an unbounded finite separation.
(b) The frozen heterogeneous fairness-preserving null triggers the homogeneous
raw-contrast predicate on two positive-measure tape intervals.  Both panels are
analytic consequences of the frozen fixtures, not Monte Carlo estimates.}
\label{fig:p17-audit}
\end{figure}


The frozen P12/P15 oracle campaign has now been independently admitted from
complete per-configuration records.  It contains 54 paired-replay and 144
four-way-replay configurations, each with 100,000 trials.  Revalidation checked
594 empirical rates against their exact binomial laws: zero comparisons crossed
the familywise $0.05$ Bonferroni threshold, the minimum exact $p$-value was
$4.6476\times10^{-4}$, and the maximum absolute empirical deviation was
$2.4271\times10^{-3}$.  The raw-result SHA-256 digests are
\texttt{f00458f5\ldots 6f1035} (P12) and
\texttt{e0750e5d\ldots 274ab} (P15).  The accompanying evidence directory
contains the full records, admission manifest, and machine-readable validation
report.  These checks validate the sufficient-statistic simulator against the
declared finite laws; they are neither proof evidence nor runtime measurements
for executable software.


\section{Experimental protocol and falsification tests}
\paragraph{Empirical design validity.}
The CPU historical-version subject and GPU sufficient-statistic simulators answer different questions. The Fairlearn table is a complete deterministic census of native-versus-canonical protected-tuple encoding discrepancies on finite evaluation rows. It is not a causal fairness estimate; therefore we do not add confidence intervals or infer demographic discrimination. Conversely, the frozen P12/P15 simulation campaigns target error behavior under theorem-specified Bernoulli families. Their admitted 198-row, 100,000-trial-per-configuration records include exact source revision, specification hashes, sampling laws, seeds, independent worker assignment, all outcomes, and merged-result hashes. A passed unit test or a completed accelerator job alone is not a measured error-rate estimate; the admitted evidence was instead checked against 594 exact finite laws.

\paragraph{Falsification-oriented validation.}
The validation cases compare the Gaussian null construction with each graph-family closed form, examine the one-edge $\tau\downarrow0$ boundary of \Cref{eq:ktau}; contrast the P17 witness-cover and raw-information geometries on a clique and an infeasible-witness graph; and distinguish fixed-allocation from sequential costs on disjoint channels. The software-version comparison holds evaluation rows, protected-group tuple ordering, admissible group support, historical dependencies, and source bytes fixed. All five exact-zero populations are retained.

\section{Conclusion}
Reusing software-evolution information can eliminate old obligations without executions, and replayable randomness can change the small-effect cost of the remaining obligations. The attainable gain depends sharply on the fairness tolerance and on which version/input combinations can share randomness. Graph allocation, partial-replay covering, and matched controls formalize that dependence. Future work should test the boundary between known and unknown controls and characterize composite-nuisance replay without overstating currently scoped guarantees.
